\section{An application: properties of lambda terms} \label{sec:Applications}
To illustrate the general applicability of our development, we show how it binds to another component of our greater project, the lambda calculus.

Our formalization of the lambda calculus is based on the so-called \emph{nested types} encoding of lambda terms.
\cite{BirdPaterson} \cite{AltenkirchReus}.
(We learned of this encoding from H\aa{}kon Gylterud; it appears to be quite popular in the Haskell community.)
The key to this representation is to think of lambda terms not as a set (type), but as a functor (a type-indexed family): 
\newcommand{\Lam}{\Lambda}
\newcommand{\tSet}{\mathtt{Set}}
\newcommand{\mtt}[1]{\mathtt{#1}}
\begin{align*}
  \mtt{data}&\ \Lam\ (X : \tSet) : \tSet \ \ \mtt{where} \ \ \\ 
  \mtt{var} &: X → \Lam X \\
  \mtt{app} &: \Lam X → \Lam X → \Lam X \\
  \mtt{abs} &: \Lam (\uparrow\! X) → \Lam X
\end{align*}
The $\uparrow$ type constructor represents the lifting monad $\uparrow X = X + 1$. 
Its occurrence inside the argument to $\Lam$ in the $\mtt{abs}$ constructor renders the type non-homogeneous: the type of lambda terms over a set $A$ 
depends on the type of terms over other sets.

The advantages of this encoding include the following:
\begin{itemize}
  \item Lifting/dropping lemmas related to de Bruijn indices obtain natural categorical forms;
  \item The encoding suggests a general higher-order induction combinator that suffices to derive all of the standard results like the substitution lemma;
    % standard facts from the classical theory;
  \item The framework is powerful enough to do formalization of classical lambda calculus theory.
\end{itemize}
\newcommand{\mcH}{\mathcal{H}}
\newcommand{\mcHo}{\mathcal{H\omega}}
The current state of the library includes standard results about reduction (standardization and confluence), 
definitions of undecidable lambda theories based on identifying unsolvable terms ($\mcH$ and $\mcHo$), 
and a formalization of an unpublished result on characterizing 2-cycles, based on earlier collaboration with Jan Willem Klop and Joerg Endrullis.

As an example of how the ARS library helps this development, the following basic implications are obtained via the ARS theory:
\begin{description}
  \item[$\SN \to \WN$:] Agda can effectively compute the normal form of any SN term.
  \item[$\CR \to \UN$:] Distinct normal forms are not beta convertible (equational consistency).
  \item[$\mathtt{dec NF_\beta}$:] It is decidable whether a given term is a normal form.
  \item[Well-foundedness:] Most properties displayed in Figure \ref{fig:WF} are equivalent for lambda terms.
\end{description}
While all these facts are easy and could have been proved directly, that they can be obtained as consequences of the general theory 
illustrates the whole point: such facts shouldn't have to be proved at all, and now they don't have to be.

An additional aspect of this application worth of note is that the heterogeneous nature of the $\Lam$ type family in no way interfered with the classical, Terese-style 
formalization of ARS. For example, the proof of decidability of normal forms is derived from the fact that single-step beta reduction is 
decidable and finitely branching, concepts that are defined for a fixed domain $A$ --- even though a single lambda abstraction requires to consider 
reduction between terms over a different set, instantiating ARS theory in a new domain.
This shows that Agda's module system indeed interacts very well with its type system. 
It was a surprise to see a design feature for namespace management to be so useful in relating results with a high level of generality.

